\documentclass[11pt]{article}

\usepackage[utf8]{inputenc}
\usepackage[T1]{fontenc}
\usepackage[spanish]{babel}
\usepackage{amsmath,amssymb,mathtools,bm,graphicx,xcolor}
\usepackage{svg}
\svgsetup{inkscapelatex=false}
\usepackage[a4paper,margin=2.5cm]{geometry}
\usepackage[
    colorlinks=true,
    linkcolor=red,
    citecolor=green,
    urlcolor=red
]{hyperref}

\spanishdecimal{.}

\setlength{\parindent}{0pt}
\setlength{\parskip}{0.45em}
\setlength{\abovedisplayskip}{6pt}
\setlength{\belowdisplayskip}{6pt}
\setlength{\abovedisplayshortskip}{4pt}
\setlength{\belowdisplayshortskip}{4pt}

% Formato de encabezado y problemas
\newcommand{\encabezado}[3]{%
{\Large\bfseries #1\par}
\vspace{2pt}
{\normalsize #2\hfill #3\par}
\vspace{6pt}\hrule\vspace{8pt}}

\newcommand{\problema}[2]{%
\vspace{0.55em}
\noindent\textbf{#1.}\enspace{\small #2}\par
\vspace{0.3em}}

\newcommand{\inciso}[1]{%
\vspace{0.35em}
\noindent\textbf{(#1)}\enspace}

% Comandos auxiliares
\newcommand{\dd}{\,\mathrm{d}}
\newcommand{\ii}{\mathrm{i}}
\newcommand{\e}{\mathrm{e}}
\newcommand{\sen}{\operatorname{sen}}
\newcommand{\grad}{\nabla}
\newcommand{\laplaciano}{\nabla^2}
\newcommand{\R}{\mathbb{R}}

\begin{document}

\encabezado{Armónicos esféricos e integral angular}{Física Matemática II}{J.R., G.R.}

\problema{1}{Muestre las expresiones cartesianas de los armónicos esféricos hasta $\ell=2$ y use esos resultados para calcular
\begin{equation}
\label{integral-armonicos}
I=\int_{0}^{\pi}\int_{0}^{2\pi}(x+y+2z)Y_\ell^m(\theta,\varphi)\sen\theta\dd\varphi\dd\theta.
\end{equation}}

\textbf{Solución:} Consideremos las relaciones usuales entre coordenadas cartesianas y coordenadas esféricas:
\begin{equation}
\label{cartesianas-esfericas}
x=r\sen\theta\cos\varphi,
\qquad
 y=r\sen\theta\sen\varphi,
\qquad
 z=r\cos\theta.
\end{equation}
De (\ref{cartesianas-esfericas}) se obtiene además
\begin{equation}
\label{relaciones-complejas-armonicos}
\frac{z}{r}=\cos\theta,
\qquad
\frac{x\pm \ii y}{r}=\sen\theta\,\e^{\pm\ii\varphi}.
\end{equation}

\inciso{a} Los primeros armónicos esféricos son
\begin{align*}
Y_0^0&=\frac{1}{\sqrt{4\pi}},\\
Y_1^0&=\sqrt{\frac{3}{4\pi}}\cos\theta,
&
Y_1^{\pm1}&=\mp\sqrt{\frac{3}{8\pi}}\sen\theta\,\e^{\pm\ii\varphi},\\
Y_2^0&=\sqrt{\frac{5}{16\pi}}(3\cos^2\theta-1),
&
Y_2^{\pm1}&=\mp\sqrt{\frac{15}{8\pi}}\sen\theta\cos\theta\,\e^{\pm\ii\varphi},\\
Y_2^{\pm2}&=\sqrt{\frac{15}{32\pi}}\sen^2\theta\,\e^{\pm2\ii\varphi}.
\end{align*}
Sustituyendo (\ref{relaciones-complejas-armonicos}) en las expresiones anteriores, se deduce directamente que
\begin{align}
\label{armonicos-cartesianos}
Y_0^0&=\frac{1}{\sqrt{4\pi}},\nonumber\\
Y_1^0&=\sqrt{\frac{3}{4\pi}}\frac{z}{r},\nonumber\\
Y_1^{\pm1}&=\mp\sqrt{\frac{3}{8\pi}}\frac{x\pm\ii y}{r},\nonumber\\
Y_2^0&=\sqrt{\frac{5}{16\pi}}\frac{2z^2-x^2-y^2}{r^2},\\
Y_2^{\pm1}&=\mp\sqrt{\frac{15}{8\pi}}\frac{(x\pm\ii y)z}{r^2},\nonumber\\
Y_2^{\pm2}&=\sqrt{\frac{15}{32\pi}}\frac{(x\pm\ii y)^2}{r^2}.\nonumber
\end{align}

\inciso{b} Para calcular la integral (\ref{integral-armonicos}), primero escribimos $x+y+2z$ como combinación lineal de armónicos con $\ell=1$. A partir de (\ref{armonicos-cartesianos}) con $\ell=1$, tenemos
\begin{align*}
Y_1^0&=\sqrt{\frac{3}{4\pi}}\frac{z}{r},\\
Y_1^1&=-\sqrt{\frac{3}{8\pi}}\frac{x+\ii y}{r},\\
Y_1^{-1}&=\sqrt{\frac{3}{8\pi}}\frac{x-\ii y}{r}.
\end{align*}
Despejando $x$, $y$ y $z$ de estas relaciones, obtenemos
\begin{align}
\label{x-y-z-en-armonicos}
x&=r\sqrt{\frac{2\pi}{3}}\left(Y_1^{-1}-Y_1^1\right),\\
y&=\ii r\sqrt{\frac{2\pi}{3}}\left(Y_1^1+Y_1^{-1}\right),\\
z&=r\sqrt{\frac{4\pi}{3}}Y_1^0.
\end{align}
Sustituyendo (\ref{x-y-z-en-armonicos}) en $x+y+2z$, encontramos
\begin{equation}
\label{xy2z-expansion}
x+y+2z
=r\sqrt{\frac{2\pi}{3}}(1+\ii)Y_1^{-1}
+r\sqrt{\frac{2\pi}{3}}(-1+\ii)Y_1^1
+2r\sqrt{\frac{4\pi}{3}}Y_1^0.
\end{equation}

Usamos ahora la relación de ortogonalidad sin conjugación,
\begin{equation}
\label{ortogonalidad-sin-conjugar-armonicos}
\int_{S^2}Y_{\ell'}^{m'}(\theta,\varphi)Y_\ell^m(\theta,\varphi)\dd\Omega
=(-1)^m\delta_{\ell\ell'}\delta_{m',-m},
\qquad
\dd\Omega=\sen\theta\dd\theta\dd\varphi.
\end{equation}
Sustituyendo (\ref{xy2z-expansion}) en (\ref{integral-armonicos}) y usando (\ref{ortogonalidad-sin-conjugar-armonicos}), notamos que sólo contribuyen los armónicos con $\ell=1$. Por lo tanto,
\begin{equation*}
I=0,
\qquad \ell\neq1.
\end{equation*}
Para $\ell=1$, evaluamos caso a caso:
\begin{equation*}
I=
\left\{
\begin{array}{lcc}
r\sqrt{\dfrac{2\pi}{3}}(1-\ii), & \text{si} & m=-1,\\[0.8em]
2r\sqrt{\dfrac{4\pi}{3}}, & \text{si} & m=0,\\[0.8em]
-r\sqrt{\dfrac{2\pi}{3}}(1+\ii), & \text{si} & m=1,\\[0.8em]
0, & \text{si} & \ell\neq1.
\end{array}
\right.
\end{equation*}
Si la integral se toma sobre la esfera unitaria, entonces basta fijar $r=1$ en la expresión anterior.



\section*{Visualización numérica}

En las figuras siguientes se representa la dependencia angular de las funciones sobre la esfera unitaria. La coordenada horizontal es la longitud $\varphi$ y la coordenada vertical es la latitud angular
\begin{equation*}
\lambda=\frac{\pi}{2}-\theta,
\end{equation*}
de modo que $\lambda=0$ corresponde al ecuador, $\lambda=\pi/2$ al polo norte y $\lambda=-\pi/2$ al polo sur. El color representa el valor de la función evaluada en cada dirección $(\theta,\varphi)$. Las líneas negras, cuando aparecen, indican curvas nodales, es decir, curvas donde la cantidad graficada se anula.

\begin{figure}[h]
\centering
\includesvg[width=0.82\textwidth]{figures/fig_01_mapa_funcion_angular}
\caption{Mapa angular firmado de $(x+y+2z)/r$ sobre la esfera unitaria. El color indica el signo y la magnitud relativa de esta función angular. La línea negra corresponde al nodo $(x+y+2z)/r=0$, que separa las regiones de signo positivo y negativo.}
\end{figure}

\begin{figure}[h]
\centering
\includesvg[width=0.82\textwidth]{figures/fig_02_mapa_magnitud_funcion_angular}
\caption{Mapa de la magnitud angular $|(x+y+2z)/r|$. En este caso el color no distingue signo, sino sólo tamaño relativo. Los máximos se ubican en las direcciones donde el vector $(1,1,2)$ queda alineado o antialineado con la normal de la esfera.}
\end{figure}

\begin{figure}[h]
\centering
\includesvg[width=0.98\textwidth]{figures/fig_03_mapas_base_l1}
\caption{Mapas angulares de tres modos reales asociados al subespacio $\ell=1$. Estos modos son proporcionales a $x/r$, $-y/r$ y $z/r$, respectivamente. El color representa el valor del modo sobre cada dirección de la esfera y las líneas negras indican los nodos donde cada modo se anula.}
\end{figure}

\begin{figure}[h]
\centering
\includesvg[width=0.88\textwidth]{figures/fig_04_mapas_l2}
\caption{Ejemplos de modos reales con $\ell=2$. Su estructura angular posee más regiones de cambio de signo que los modos con $\ell=1$. Por esta diferencia de simetría, estos modos no contribuyen al proyectar una función lineal en $x$, $y$ y $z$ sobre la base de armónicos esféricos.}
\end{figure}

\begin{figure}[h]
\centering
\includesvg[width=0.98\textwidth]{figures/fig_05_densidades_integracion}
\caption{Mapas angulares de productos que aparecen bajo la integral de proyección. El color muestra qué regiones de la esfera aportan positiva o negativamente a la integral. Las líneas negras indican nodos del integrando. En los productos compatibles con $\ell=1$ queda una contribución neta, mientras que la contribución con un modo $\ell=2$ se cancela por ortogonalidad angular.}
\end{figure}

\end{document}
